\section{Time-dependent Green's functions}

\subsection{Wave equation}

Consider the forced vibration problem
\begin{equation}
\label{eq:wave-forced}
u_{tt}-a^2\nabla^2 u=f(\mathbf{r},t),
\end{equation}
\begin{equation}
\Bigl[\alpha\frac{\partial u}{\partial n}+\beta u\Bigr]_{\Sigma}=\theta(M,t),
\end{equation}
\begin{equation}
u\big|_{t=0}=\varphi(\mathbf{r}),
\qquad
u_t\big|_{t=0}=\psi(\mathbf{r}).
\end{equation}
A continuous force density is a superposition of impulses,
\begin{equation}
f(\mathbf{r},t)
=\iiint_{T}\int_{-\infty}^{\infty}
f(\mathbf{r}_0,\tau)\,
\delta(\mathbf{r}-\mathbf{r}_0)\,\delta(t-\tau)\,
\mathrm{d}\mathbf{r}_0\,\mathrm{d}\tau.
\end{equation}
The \textbf{Green's function} $G(\mathbf{r},t;\mathbf{r}_0,t_0)$ is the response to a unit impulse at $(\mathbf{r}_0,t_0)$ with $t_0>0$:
\begin{equation}
\label{eq:wave-G}
G_{tt}-a^2\nabla^2 G=\delta(\mathbf{r}-\mathbf{r}_0)\delta(t-t_0),
\end{equation}
\begin{equation}
\Bigl[\alpha\frac{\partial G}{\partial n}+\beta G\Bigr]_{\Sigma}=0,
\end{equation}
\begin{equation}
G\big|_{t=0}=0,
\qquad
G_t\big|_{t=0}=0.
\end{equation}

\subsubsection{Reciprocity with time reversal}

The wave Green's function satisfies
\begin{equation}
\label{eq:wave-recip}
G(\mathbf{r},t;\mathbf{r}_0,t_0)=G(\mathbf{r}_0,-t_0;\mathbf{r},-t).
\end{equation}
The idea of the proof is to introduce the time-reversed field $G(\mathbf{r},-t;\mathbf{r}_1,-t_1)$, multiply and subtract against \eqref{eq:wave-G}, integrate over $T\times(-\infty,t']$, convert the Laplacian terms by Green's second identity, and use the homogeneous boundary conditions together with causality.

\subsubsection{Integral representation}

Writing the original problem in the source variables $(\mathbf{r}_0,t_0)$ and combining it with the reciprocal form of $G$ (causal: $G=0$ for $t<t_0$), one multiplies, subtracts, and integrates over $T\times[0,t+\varepsilon]$. Using the identity
\begin{equation}
G u_{t_0t_0}-u G_{t_0t_0}
=\partial_{t_0}\bigl(G u_{t_0}-u G_{t_0}\bigr)
\end{equation}
and Green's identity in space, one obtains
\begin{equation}
\label{eq:wave-rep}
\begin{aligned}
u(\mathbf{r},t)
&=\iiint_{T}\int_0^{t}G(\mathbf{r},t;\mathbf{r}_0,t_0)f(\mathbf{r}_0,t_0)\,\mathrm{d}V_0\,\mathrm{d}t_0\\
&\quad+a^2\iint_{\Sigma}\int_0^{t}
\Bigl(G\frac{\partial u}{\partial n_0}-u\frac{\partial G}{\partial n_0}\Bigr)
\,\mathrm{d}S_0\,\mathrm{d}t_0\\
&\quad+\iiint_{T}\bigl[G u_{t_0}-u G_{t_0}\bigr]_{t_0=0}\,\mathrm{d}V_0.
\end{aligned}
\end{equation}
The three terms are, respectively, the accumulated effect of the forcing, the contribution of the boundary data, and the imprint of the initial conditions.

\subsection{Heat (transport) equation}

For the inhomogeneous heat problem
\begin{equation}
u_t-a^2\nabla^2 u=f(\mathbf{r},t),
\end{equation}
\begin{equation}
\Bigl[\alpha\frac{\partial u}{\partial n}+\beta u\Bigr]_{\Sigma}=\theta(M,t),
\end{equation}
\begin{equation}
u\big|_{t=0}=\varphi(\mathbf{r}),
\end{equation}
an analogous calculation yields
\begin{equation}
\label{eq:heat-rep}
\begin{aligned}
u(\mathbf{r},t)
&=\iiint_{T}\int_0^{t}Gf\,\mathrm{d}V_0\,\mathrm{d}t_0\\
&\quad+a^2\iint_{\Sigma}\int_0^{t}
\Bigl(G\frac{\partial u}{\partial n_0}-u\frac{\partial G}{\partial n_0}\Bigr)
\,\mathrm{d}S_0\,\mathrm{d}t_0\\
&\quad+\iiint_{T}\bigl[uG\bigr]_{t_0=0}\,\mathrm{d}V_0.
\end{aligned}
\end{equation}
Only a single initial datum appears, because the heat equation is first order in time.

\begin{note}
As in the elliptic case, one chooses $G$ to satisfy the homogeneous boundary condition of the same type as the given data, so that the unknown combination of $u$ and $\partial u/\partial n$ on $\Sigma$ drops out of \eqref{eq:wave-rep} or \eqref{eq:heat-rep}.
\end{note}
